\documentclass[compress,orange,subsection=false,10pt]{beamer}
\usepackage[latin1]{inputenc}
\usepackage[spanish]{babel}
\usepackage{verbatim}
\usepackage{listings}
\usepackage{subfigure}
\usepackage{multicol}
\usepackage{amsmath}
\usepackage{epsfig}
\usepackage{graphicx}
\usepackage[all,knot]{xy}
\xyoption{arc}
\usepackage{url}
\usepackage{multimedia}
\usepackage{hyperref}

\mode<presentation>
\usetheme{Madrid}
%\usecolortheme[RGB={111,73,135}]{structure}
\usecolortheme[RGB={128,0,0}]{structure}
%\usecolortheme[RGB={0,96,0}]{structure}
%\usecolortheme[RGB={200,0,200}]{structure}
%\usecolortheme[RGB={0,128,0}]{structure}
%\usecolortheme[RGB={0,0,128}]{structure}
\usefonttheme{serif}
\useinnertheme{rectangles}
\useoutertheme{split}
\setbeamercovered{transparent}

\title{\LaTeX{} Beamer}
\subtitle{Presentaciones como dios manda}

\author{Jes�s Espino Garc�a}
\institute{\includegraphics[height=1.5cm]{gul}}

\date{10 de Octubre de 2008}

%\logo{\includegraphics[height=1.5cm]{gul}}

\xdefinecolor{olive}{cmyk}{0.64,0,0.95,0.4}
\xdefinecolor{purpleish}{cmyk}{0.75,0.75,0,0}

\begin{document}
\frame{\titlepage}

\begin{frame}{�Que es?}
	\begin{itemize}
		\item Extension de \LaTeX{}.
		\item A�ade una clase de documento para transparencias.
		\item Desarrollado por Till Tantau desde el 2003.
		\item Actualmente esta en la version 3.
	\end{itemize}
\end{frame}

\begin{frame}{�Por Que?}
	\transblindsvertical<1>
	\only<1>{
		\begin{block}{Pros}
			\begin{itemize}
				\item Es \LaTeX{}!!
				\item WYSIWYM.
				\item Efectos y transiciones.
				\item Aspecto adaptable.
				\item Plantillas.
				\item Salidas en diferentes formatos. 
				\item Barras de navegacion.
			\end{itemize}
		\end{block}
	}
	\only<2>{
		\begin{block}{Contras}
			\begin{itemize}
				\item Es \LaTeX{} :(
				\item WYSIWYM.
				\item Peque�o esfuerzo inicial.
				\item Adaptaciones profundas del aspecto complicadas.
			\end{itemize}
		\end{block}
	}
\end{frame}


\section{Elementos}
\subsection{Preludio}
\begin{frame}[fragile]{El preludio}
	\begin{block}{Preludio}
		\begin{enumerate}
			\item \verb+\documentclass{beamer}+
			\item \verb+\usetheme{Warsaw}+
			\item \verb+\author{Jes�s Espino}+
			\item \verb+\title{\LaTeX Beamer}+
			\item \verb+\date{10 de Octubre de 2008}+
			\item \verb+\begin{document}+
		\end{enumerate}
	\end{block}
\end{frame}

\subsection{Frames}

\begin{frame}{Los frames}
	\transdissolve<1-3>
	\begin{block}{Sobre los frames}
		\begin{itemize}[<+->]
			\item Es un entorno de \LaTeX{}.
			\item Cada frame es un transparencia.
			\item Cada frame puede tener un titulo y un subtitulo.
			\item Cada frame tiene unas opciones.
		\end{itemize}
	\end{block}
\end{frame}

\begin{frame}[containsverbatim]{Como se escribe}
	\begin{block}{Frame}
		\begin{verbatim}
\begin{frame}[opciones]{Titulo}{SubTitulo}
   Contenido de la transparencia
\end{frame}
		\end{verbatim}
	\end{block}
\end{frame}

\begin{frame}{Opciones de los frames}
	\begin{block}{Opciones}
	\begin{columns}
		\column{0.45\textwidth}
		\begin{itemize}
			\item allowframebreak
			\item allowdisplaybreaks
			\item b
			\item c
			\item t
			\item fragile
		\end{itemize}
		\column{0.45\textwidth}
		\begin{itemize}
			\item environment
			\item label
			\item plain
			\item shrink
			\item squeeze
		\end{itemize}
	\end{columns}
	\end{block}
\end{frame}

\begin{frame}[containsverbatim]{Ejemplo de frame}
	\begin{block}{Ejemplo}
		\begin{verbatim}
\begin{frame}{Ejemplo de frame}
    \begin{itemize}
        \item Esto es un ejemplo.
        \item Tiene items.
        \item \verb+Algunos en modo verbatim+.
    \end{itemize}
\end{frame}
		\end{verbatim}
	\end{block}
\end{frame}

\subsection{Bloques}

\begin{frame}[containsverbatim]{Las bloques}
	Nos permiten resaltar y agrupar contenidos:
	\begin{block}{Ejemplo de bloque con contenidos}
		$\left|{1\over N}\sum_{n=1}^N \gamma(u_n)\right|$
		\begin{itemize}
			\item Una formula y un listado.
			\item Dentro de un bloque.
		\end{itemize}
	\end{block}
	\begin{block}{Codigo del bloque}
		\begin{verbatim}
\begin{block}{Ejemplo de bloque con contenidos}
  $\left|{1\over N}\sum_{n=1}^N \gamma(u_n)\right|$
  \begin{itemize}
    \item Una formula y un listado.
    \item Dentro de un bloque.
  \end{itemize}
\end{block}
		\end{verbatim}
	\end{block}
\end{frame}

\subsection{Columnas}

\begin{frame}{Las columnas}
	Nos permite dividir en varias columnas nuestros contenidos:
	\begin{columns}[c]
		\column{0.5\textwidth}
			\begin{block}{Izquierda}
				\begin{itemize}
					\item Un listado
					\item Dentro de un bloque
					\item Dentro de una columna
				\end{itemize}
			\end{block}
		\column{0.5\textwidth}
			\begin{block}{Derecha}
				$\left|{1\over N}\sum_{n=1}^N \gamma(u_n)-{1\over 2\pi}\int_0^{2\pi}\gamma(t){\rm d}t\right|$
				\begin{itemize}
					\item Una formula y un listado.
					\item Dentro de un bloque.
					\item Dentro de una columna.
				\end{itemize}
			\end{block}
	\end{columns}
\end{frame}

\begin{frame}{Opciones de las columnas}
	\begin{itemize}
		\item b: Alinea las columnas por la ultima linea.
		\item c: Alinea las columnas por el centro.
		\item t: Alinea las columnas por la primera linea (por defecto).
		\item T: Alinea las columnas por su punto mas arriba.
		\item totalwidth: Tama�o maximo del conjunto de columnas.
		\item onlytextwidth igual que totalwidth=\\textwidth.
	\end{itemize}
\end{frame}

\begin{frame}[containsverbatim]{Ejemplo de columnas}
	\begin{block}{Ejemplo}
		\begin{verbatim}
\begin{columns}[c]
  \column{0.5\textwidth}
    \begin{block}{Izquierda}
      \begin{itemize}
        \item Un listado
      \end{itemize}
    \end{block}
  \column{0.5\textwidth}
  \begin{block}{Derecha}
    $\left|{1\over N}\sum_{n=1}^N \gamma(u_n)...$
    \begin{itemize}
      \item Una formula y un listado.
    \end{itemize}
  \end{block}
\end{columns}
		\end{verbatim}
	\end{block}
\end{frame}

\section{Aspecto}
\subsection{Tema}
\begin{frame}[containsverbatim]{Tema}
	Nos permite configurar todo el layout de nuestra presentacion
	\begin{block}{Uso}
		\verb+\usetheme[ options ]{ name list }+
	\end{block}
	\begin{block}{Ejemplo}
		\verb+\usetheme[right]{Berkeley}+
	\end{block}
\end{frame}

\subsection{Color}
\begin{frame}[containsverbatim]{Color}
	Nos permite configurar el conjunto de colores del layout de nuestra presentacion
	\begin{block}{Uso}
		\verb+\usecolortheme[ options ]{ name list }+
	\end{block}
	\begin{block}{Ejemplo}
		\verb+\usecolortheme[ RGB={128,0,0} ]{ structure }+
	\end{block}
\end{frame}

\subsection{Fuente}
\begin{frame}[containsverbatim]{Fuente}
	Nos permite configurar la fuente del layout de nuestra presentacion
	\begin{block}{Uso}
		\verb+\usefonttheme[ options ]{ name }+
	\end{block}
	\begin{block}{Ejemplo}
		\verb+\usefonttheme{ serif }+
	\end{block}
\end{frame}

\subsection{Interior}
\begin{frame}[containsverbatim]{Interior}
	Nos permite configurar la parte interior del layout de nuestra presentacion
	\begin{block}{Uso}
		\verb+\useinnertheme[ options ]{ name }+
	\end{block}
	\begin{block}{Ejemplo}
		\verb+\useinnertheme{ rectangles }+
	\end{block}
\end{frame}

\subsection{Exterior}
\begin{frame}[containsverbatim]{Exterior}
	Nos permite configurar la parte exterior del layout de nuestra presentacion
	\begin{block}{Uso}
		\verb+\useoutertheme[ options ]{ name }+
	\end{block}
	\begin{block}{Ejemplo}
		\verb+\useoutertheme[ footline=authortitle ]{ miniframes }+
	\end{block}
\end{frame}


\section{Efectos}
\subsection{Overlays}
\begin{frame}[containsverbatim]{Overlays}
	\begin{itemize}
		\item Separan nuestra presentacion en capas.
		\item Estas capas puede ocultarse o mostrarse.
		\item Se les puede establecer un orden de aparicion.
		\item Se suelen especificar con la expresion \verb+<n>+.
		\item Ej: \verb+\item<1>+, \verb+\begin{block}{Bloque}<1-3>\end{block}+, \verb+\item<3->+.
		\item Se pueden establecer pausas con \verb+\pause+.
		\item Se pueden ocultar o omitir los contenidos \verb+\only<1> != \item<1>+
	\end{itemize}
\end{frame}

\begin{frame}{Ejemplos de Overlays}
	\begin{block}{Bloque1}<1>Este es el bloque 1\end{block}
	\begin{block}{Bloque2}<2>Este es el bloque 2\end{block}
	\only<3>{
		\begin{block}{Bloque 3}Este es el bloque 3\end{block}
	}
	\begin{block}{Bloque4}<1->\color<2-4>[rgb]{1,0,0}Este es el bloque 4\end{block}
\end{frame}

\subsection{Transiciones}
\begin{frame}{Transiciones}
	\transblindsvertical<1>
	\transboxin<2>
	\transboxout<3>
	\transdissolve<4>
	\transglitter<5>
	\begin{block}{transblindsvertical}<1>\end{block}
	\begin{block}{transboxin}<2>\end{block}
	\begin{block}{transboxout}<3>\end{block}
	\begin{block}{transdissolve}<4>\end{block}
	\begin{block}{transglitter}<5>\end{block}
\end{frame}
\begin{frame}{Transiciones}
	\transsplitverticalin<1>
	\transsplitverticalout<2>
	\transsplithorizontalin<3>
	\transsplithorizontalout<4>
	\transwipe<5>
	\begin{block}{transsplitverticalin}<1>\end{block}
	\begin{block}{transsplitverticalout}<2>\end{block}
	\begin{block}{transsplithorizontalin}<3>\end{block}
	\begin{block}{transsplithorizontalout}<4>\end{block}
	\begin{block}{transwipe}<5>\end{block}
\end{frame}

\section{Referencias}
\begin{frame}{Referencias}
	\begin{block}{Referencias}
		\begin{itemize}
			\item Latex Beamer User Guide
			\item Lista de correo del GUL <gul@gul.uc3m.es>
		\end{itemize}
	\end{block}
\end{frame}

\section{Dudas}
\begin{frame}{Dudas}
	\dots{}
\end{frame}
\end{document}
